\chapter{Décompositions tensorielles}

Les données pour le TD sont stockées dans le fichier \lstinline{td_tenseurs.mat} (elles peuvent être chargées avec \lstinline{load}).



Pour cet excercice, il faut utiliser le \emph{N-way Toolbox} qui peut être trouvé dans le fichier {\tt nway331.zip}.
Il est mieux de mettre le toolbox dans un dossier séparé (par exemple {\tt nway}) et utiliser {\tt addpath}. 

\section{Données de spectroscopie}

Les variables {$\mathbf{A}, \mathbf{B}, \mathbf{C}$ stockés dans \lstinline{td_tenseurs.mat}  sont :
\[
\begin{split}
\mathbf{A} = \begin{bmatrix} \mathbf{a}_1 & \cdots & \mathbf{a}_R \end{bmatrix} \in \mathbb{R}^{I \times R},\;
\mathbf{B} = \begin{bmatrix} \mathbf{b}_1 & \cdots & \mathbf{b}_R \end{bmatrix} \in \mathbb{R}^{J \times R},\;
\mathbf{C} = \begin{bmatrix} \mathbf{c}_1 & \cdots & \mathbf{c}_R \end{bmatrix} \in \mathbb{R}^{K \times R}.
\end{split}
\]
Ce sont les données issues de spectroscopie de fluorescence (mélange de $R$ composants, $\mathbf{a}_k$ --- concentrations, $\mathbf{c}_k$ --- spectres d'excitation; $\mathbf{b}_k$ --- spectres d'émission).

\begin{enumerate}
\item Tracer les concentrations, spectres d'excitations et spectres d'émissions à l'aide de {\tt plot}.
\item Construire le tenseur $\mathcal{T} \in \mathbb{R}^{I \times J \times K}$ donné par 
\[
\mathcal{T} = \sum_{k=1}^{R} \mathbf{a}_k \circ  \mathbf{b}_k  \circ  \mathbf{c}_k,
\]
comme dans l'exercice précédent.

\item Sélectionner la tranche horizontale $\mathbf{Z} = \mathcal{T}_{3,:,:}$. Afficher la tranche avec {\tt imagesc}.

\end{enumerate}

\section{Décompositions}
\begin{enumerate}
\item Verifier  que la matrice $\mathbf{Z}$ admet une decomposition
\[
\mathbf{Z} = \mathbf{B} \begin{bmatrix} A_{3,1} & 0 & 0 \\
0 & \ddots & 0 \\ 0&0& A_{3,R}\end{bmatrix}   \mathbf{C}^{\top}.
\]
La matrice diagonale peut être construite à l'aide de la fonction {\tt diag} à partir de la troisième ligne de $\mathbf{A}$.

\item Calculer une décomposition  CP   de rang $R$ avec {\tt parafac(T,R)}. Peut-on trouver les vecteurs $\mathbf{a}_k, \mathbf{b}_k,\mathbf{c}_k$ de la décomposition CP initiale? 

%\item * Peut-on retrouver les spectres à partir de $\mathbf{Z}$ avec la NMF?  Les conditions d'unicité (page 24 dans {\tt NMF\_ISC.pdf}) sont-elles satisfaites?

\item Calculer la SVD tronquée  
\[
\mathbf{Z} =  \mathbf{U}_{:,1:R} \mathbf{\Sigma}_{1:R,1:R} \mathbf{V}_{:,1:R}^{\top}.
\]
 Peut-on retrouver les spectres à partir des vecteurs singuliers? (Tracez les $R$ premiers vecteurs singuliers droits et  comparez-les avec  les $\mathbf{c}_k$.) 
\end{enumerate}




